\worldchapter{Spielkonzept}{game-concept}
\label{ch:game-concept}

\begin{multicols*}{2}

Obwohl die Spielleiterin die Bösen steuert, steht sie eigentlich auf der Seite der Helden.
Spieler und Spielleiterin begeben sich zusammen auf eine fantastische Reise, bei der die Erzählung im Vordergrund steht.

Die Spieler steuern ihre \textbf{Charaktere} in dieser fiktiven Welt, und handeln in ihrem Namen.
Die Fähigkeiten eines Charakters, ihr Hab und Gut, ihre Stärken und Schwächen werden auf dem \textbf{Charakterbogen} niedergeschrieben.
Du als Spieler beschreibst der Spielleiterin das Handeln deines Charakters.
Möglicherweise erfordert die Handlung deines Charakters eine \textbf{Probe}, einen Test auf eine bestimmte Eigenschaft oder eine Fertigkeit.

Diese Proben sind der Kern des Spiels.
Es ist die Grundmechanik, um über den Erfolg oder das Misslingen einer Handlung zu bestimmen.

\section{Die Probe}

Die Probe\index{Probe} ist der zentrale Bestandteil der Regeln.
Wenn eine Situation einen ungewissen Ausgang hat, wird eine Probe verwendet, um den Ausgang der Situation zu bestimmen.
Tirakans Reiche benutzt für alle Proben einen hundertseitigen Würfel, einen \textbf{W100}.
In der Regel werden zwei zehnseitige Würfel (\textbf{W10}) geworfen, wobei einer der Würfel die Zehnerstelle und der andere die Einerstelle des Wurfes darstellen.
Ein Ergebnis von \textbf{00} bedeutet 100.

Ob Erkundung, soziale Spannung, Verfolgung, Ritual oder Kampf: Wenn der Ausgang unsicher ist und Konsequenzen hat, wird eine Probe abgelegt.

\subsection{Zielwert}

Eine Probe wird immer auf einen Wert geworfen, der sich aus einem \textbf{Attribut}\index{Attribut} und einer \textbf{Fertigkeit}\index{Fertigkeit} zusammen setzt.
Attribut und Fertigkeit werden in der jeweiligen Situation von der Spielleiterin vorgegeben, wenn die Regeln nicht etwas anderes vorsehen.

Dieser \textbf{Zielwert}\index{Zielwert} muss unterschritten werden, damit die Probe erfolgreich ist.
Der Zielwert ergibt sich, indem Attribut und Fertigkeit mit 10 multipliziert und dann addiert werden, und 30 hinzugefügt werden.

\[
Zielwert = 30 + (Attribut \times 10) + (Fertigkeit \times 10)
\]

Der Zielwert wird immer auf 5 bis 95 begrenzt.
Eine Probe gelingt, wenn der Wurf kleiner als der Zielwert ist.

\section{Schwierigkeit}

Eine Probe kann erschwert sein, wenn sie unter schwierigen Bedingungen ausgeführt wird.
Die Spielleiterin legt hier einen Malus fest, der vor dem Wurf auf den Zielwert angerechnet wird, bevor geworfen wird.

\begin{itemize}[leftmargin=*]
\item \textbf{Leicht}: -20
\item \textbf{Standard}: +0
\item \textbf{Schwer}: -20
\item \textbf{Extrem}: -40
\end{itemize}

\section{Ablauf einer Probe}

\begin{enumerate}[leftmargin=*]
\item Der Spieler beschreibt Absicht und Vorgehen.
\item Die Spielleiterin bestimmt das passende Attribut und die passende Fertigkeit.
\item Die Spielleiterin nennt Lage und Effekt, falls die Probe wichtig ist.
\item Alle Modifikatoren werden zusammengerechnet.
\item Es wird 1W100 gewürfelt.
\item Ergebnisstufe, unmittelbare Wirkung und Konsequenz werden sofort abgehandelt.
\end{enumerate}

\section{Ergebnisstufe}

Das Resultat der Probe ergibt die Ergebnisstufe.
Ist der geworfene Wert bis zu 10 Punkte vom Zielwert entfernt, so ist die Probe knapp gelungen oder misslungen.
Weicht er mehr als 30 Punkte ab, so handelt es sich um einen starken Erfolg oder schweren Fehlschlag.
Ein Wurf von 01--05 ist immer ein kritischer Erfolg, alles von 96--100 ist immer ein katastrophales Scheitern.

\begin{tabular}{ll}
    \toprule
    \textbf{Ergebnisstufe}            & \textbf{Wurf}          \\
    \midrule
    \textbf{Kritischer Erfolg}        & 01--05   \\
    \textbf{Starker Erfolg}           & Marge 30+        \\
    \textbf{Sauberer Erfolg}          & Marge 10--29     \\
    \textbf{Knapp geschafft}          & Marge 0--9       \\
    \textbf{Knapp gescheitert}        & Marge 1--9   \\
    \textbf{Klar gescheitert}         & Marge 10--29 \\
    \textbf{Schwer gescheitert}       & Marge 30+    \\
    \textbf{Katastrophal gescheitert} & 96--00    \\
    \bottomrule
\end{tabular}

\section{Konsequenzen}

Jede misslungene Probe kann zu \textbf{Konsequenzen}\index{Konsequenz} führen, die über die eigentliche Intention der Probe hinaus gehen.
Konsequenzen werden in drei Kategorien eingeteilt und haben unterschiedliche Folgen.
Diese werden von der Spielleiterin festgelegt.

%TODO: Druck?
\begin{itemize}[leftmargin=*]
\item \textbf{Gering}: kurze Verzögerung, kleiner Verlust, +1 Druck, kurze Irritation.
\item \textbf{Ernst}: Wunde, deutlicher Verlust, Eskalation, schlechtere Position.
\item \textbf{Kritisch}: schwere Wunde, Zielverschiebung, massiver Kollateralschaden, bleibende Folge.
\end{itemize}

In der Regel führt eine bestimmte Ergebnisstufe immer zu einer bestimmten Klasse an Konsequenzen, dies kann natürlich je nach Situation von der Spielleiterin verändert werden.

\begin{itemize}[leftmargin=*]
\item Knapp gescheitert \textrightarrow{} Gering
\item Klar gescheitert \textrightarrow{} Ernst
\item Schwer oder katastrophal gescheitert \textrightarrow{} Kritisch
\item Knapp geschafft in einer verzweifelten Lage \textrightarrow{} Erfolg mit geringer Konsequenz
\end{itemize}

\section{Gegenproben}

Wenn zwei Seiten direkt gegeneinander handeln, würfeln beide gegen ihren eigenen Zielwert.

\begin{itemize}[leftmargin=*]
\item Eine Seite gelingt, die andere scheitert: Die erfolgreiche Seite gewinnt.
\item Beide gelingen: Die höhere Erfolgs-Marge gewinnt.
\item Beide scheitern mit 10+ Differenz: die geringere Fehlmarge bekommt die \textbf{Oberhand}\index{Oberhand}.
\item Beide scheitern mit 0--9 Differenz: Stillstand, und die Situation verschlechtert sich fiktional.
\end{itemize}

Oberhand erlaubt genau eines:

\begin{itemize}[leftmargin=*]
\item Initiative behalten.
\item Eine Bewegung um ein Distanzsegment erzwingen.
\item Eine geringe Konsequenz verursachen.
\end{itemize}

\section{Reihenfolge bei Gleichstand}

Wenn die Abfolge relevant wird, gilt:

\begin{enumerate}[leftmargin=*]
\item höhere Ergebnisstufe
\item höhere Erfolgs-Marge
\item höherer \textbf{Instinkt}
\item höheres \textbf{Geschick}
\item niedrigerer Rohwurf
\item sonst gleichzeitig
\end{enumerate}

\section{Omen}

\textbf{Omen}\index{Omen} ist eine gemeinsame Ressource der Gruppe.
Das Maximum an Omen der Gruppe entspricht 2 plus der Hälfte des Glaubensniveaus (abgerundet).
Zu Sitzungsbeginn wird Omen auf sein Maximum aufgefüllt.

Ein Charakter oder die Gruppe kann 1 Omen ausgeben für:

\begin{itemize}[leftmargin=*]
\item +20 auf den Zielwert vor einem Wurf
\item Neu würfeln und das neue Ergebnis behalten
\item Eine Konsequenz um eine Stufe senken
\item Einmal die Reihenfolge unterbrechen und sofort handeln
\item Eine Szenenfrage stellen und eine klare Wahrheit erhalten
\end{itemize}

Die Szenenfrage darf gestellt werden, wenn der Charakter die Antwort plausibel bemerken, erinnern, deuten oder aus der Lage lesen konnte.
Die Spielleiterin entscheidet, ob die Antwort direkt, eng oder symbolisch ausfällt, aber sie muss wahr sein.

+1 Omen wird gewonnen, wenn eine große Komplikation aus Schwäche, Eid, Schuld oder Verwundbarkeit akzeptiert wird oder wenn eine Bindung mit wirklichem Preis ausgetragen wird.
\end{multicols*}
